\documentclass[10pt,a4paper]{amsart}
\usepackage[margin=19mm]{geometry}
\usepackage{amsmath,amssymb,mathtools}
\usepackage[scr=rsfs,cal=euler]{mathalfa}
\usepackage{xfrac}
\usepackage[unicode,hidelinks,bookmarks=false]{hyperref}
\newcommand{\ie}{i.e.\ }
\newcommand{\cf}{cf.\ }
\newcommand{\resp}{respectively\ }
\newcommand{\Subsection}{\subsection}
\providecommand{\nobreakdash}{\nobreak\hbox{-}}
\setlength{\emergencystretch}{2em}
\multlinegap0pt
\usepackage{bm}
\usepackage{bbm}
\usepackage{dsfont}
\usepackage{xspace}
\usepackage{cancel}
\usepackage{mathtools}
\usepackage{amsthm}
\makeatletter
\@ifundefined{theorem}{\newtheorem{theorem}{Theorem}}{}
\makeatother
\providecommand{\Eb}{\mathbb E}
\providecommand{\Pb}{\mathbb P}
\providecommand{\Var}{\operatorname{Var}}
\providecommand{\calI}{\mathcal I}
\providecommand{\calW}{\mathcal W}
\providecommand{\op}{\operatorname{op}}
\providecommand{\vecl}{\operatorname{vec}_{<}}
\title[Minimax Synthesis of Network Mechanisms]{Minimax Synthesis of Network Mechanisms}
\author{Marios Papamichalis, Regina Ruane}
\date{Source: arXiv:2606.15836v1 (CC BY 4.0)}
\begin{document}
\maketitle

\noindent\emph{Source and licence.} Minimax Synthesis of Network Mechanisms. Version arXiv:2606.15836v1 (CC BY 4.0), distributed under the Creative Commons Attribution 4.0 International licence, as recorded in the arXiv licence metadata for this exact version and shown on the abstract page https://arxiv.org/abs/2606.15836v1. Full rights notice: licenses/arxiv-2606.15836-CC-BY-4.0.txt.\par\smallskip

\noindent\textit{Editorial scope.} Counted unit: the Theorem whose source label is \texttt{thm:fixedK-noisyOR-final2} in the article (source0.tex lines 4418--4481, source label \texttt{thm:fixedK-noisyOR-final2}), delivered together with the article's own complete original proof of it (source0.tex lines 4483--4687, one \texttt{proof} environment of 205 lines). No mathematics is written, repaired or re-derived: statement and proof are the article's own text line for line. Required context retained uncounted: none. Every definition, display and label of those lines is retained unchanged. Omitted from this package and not redistributed: 1--4417, 4482--4482, 4688--6963. Nothing inside a retained excerpt is omitted or altered: every retained body is delivered exactly as the article's lines are written, so no omission receipt of any kind accompanies this package. Citations: the retained excerpts carry no citation command at all, so no bibliography entry of the article is transcribed and no reference is redistributed. Reference adaptation: none was needed and none was made; every reference inside the delivered unit resolves inside it, so no unresolved reference or citation is produced. Printed slips kept literal: no printed word, symbol, hypothesis or number of the counted statement or of its proof was altered. Layout and class: the article's own class is \texttt{article} with packages including geometry, amssymb, natbib, graphicx, setspace, hyperref, amsmath, amsthm, mathtools, bm, booktabs, multirow, array, xcolor; the wrapper is the lane's pinned \texttt{amsart} editorial wrapper at 10pt on A4, which restates the theorem-like environments the retained text needs and adds the article's own macro definitions that the retained text needs (\texttt{\textbackslash Eb}, \texttt{\textbackslash Pb}, \texttt{\textbackslash Var}, \texttt{\textbackslash calI}, \texttt{\textbackslash calW}, \texttt{\textbackslash op}, \texttt{\textbackslash vecl}), with extra wrapper packages (bm, bbm, dsfont, xspace, cancel, mathtools). The delivered numbering is the unit's own. Assets: no retained span contains a figure environment, a graphics-inclusion command, a PDF-page inclusion, a table or a TikZ picture; no image file is copied into this package, the package declares no asset dependency and no image file is redistributed with it. Licence: version arXiv:2606.15836v1 of this article is distributed under the Creative Commons Attribution 4.0 International licence, as recorded in the arXiv licence metadata for this exact version and shown on the abstract page https://arxiv.org/abs/2606.15836v1. A full rights notice is delivered at licenses/arxiv-2606.15836-CC-BY-4.0.txt. Characters: every retained excerpt is pure ASCII, so the delivered engine, XeLaTeX with its default Latin Modern text font, reports no missing character.
\begin{theorem}[Fixed-$K$ noisy-OR weights]
\label{thm:fixedK-noisyOR-final2}
Let \(K\) be fixed and
\[
  P_{ij}(w)
  =
  1-\prod_{k=1}^K(1-w_kG_{k,ij}),
  \qquad
  w\in\calW=[\underline w,\overline w]^K\subset(0,1)^K.
\]
Assume
\[
  0\le G_{k,ij}\le C\rho,\qquad
  c\rho\le P_{ij}(w)\le C\rho
\]
uniformly on \(\calW\), and
\[
  \lambda_{\min}(M^\top M)\ge c\gamma n^2\rho^2,
  \qquad
  M=[\vecl(G_1),\ldots,\vecl(G_K)].
\]
Assume further
\[
  \rho=o(\gamma),\qquad
  \gamma^3 n^2\rho\to\infty .
\]
Let \(\widehat w\) be the MLE.  Then \(\widehat w\to_p w_0\), and
\[
  \widehat w-w_0
  =
  \calI_n(w_0)^{-1}S_n(w_0)
  +o_p\{(\gamma n^2\rho)^{-1/2}\},
\]
where
\[
  S_n(w_0)
  =
  \sum_{i<j}
  \dot P_{ij}(w_0)
  \frac{A_{ij}-P_{ij}(w_0)}
       {P_{ij}(w_0)\{1-P_{ij}(w_0)\}},
\]
\[
  \dot P_{ij,k}(w)
  =
  G_{k,ij}\prod_{\ell\ne k}(1-w_\ell G_{\ell,ij}),
\]
and
\[
  \calI_n(w)
  =
  \sum_{i<j}
  \frac{\dot P_{ij}(w)\dot P_{ij}(w)^\top}
       {P_{ij}(w)\{1-P_{ij}(w)\}} .
\]
Consequently,
\[
  \|\widehat w-w_0\|_2
  =
  O_p\!\left(\frac1{\sqrt\gamma\,n\sqrt\rho}\right),
  \qquad
  \calI_n(w_0)^{1/2}(\widehat w-w_0)\Rightarrow N(0,I_K).
\]
\end{theorem}

\begin{proof}
For fixed \(K\), the noisy-OR derivatives satisfy, uniformly on
\(\calW\),
\[
  |\partial^\alpha P_{ij}(w)|\le C_\alpha \rho^{|\alpha|}
  \qquad(|\alpha|=1,2,3).
\]
In particular,
\[
  \dot P_{ij,k}(w)=G_{k,ij}+O(\rho^2).
\]
Let \(\dot M(w)\) have columns \(\vecl\{\dot P_k(w)\}\).  Then
\[
  \|\dot M(w)-M\|_{\op}\le Cn\rho^2
  =
  o(\sqrt\gamma\,n\rho),
\]
because \(\rho=o(\gamma)\).  Therefore
\[
  \sigma_{\min}\{\dot M(w)\}
  \ge
  \tfrac12\sigma_{\min}(M)
  \ge
  c\sqrt\gamma\,n\rho
\]
uniformly on \(\calW\).  Thus, for \(w_t=w_0+t(w-w_0)\),
\[
  P(w)-P(w_0)
  =
  \left\{\int_0^1\dot M(w_t)\,dt\right\}(w-w_0),
\]
and
\[
  \|P(w)-P(w_0)\|_F^2
  \ge
  c\gamma n^2\rho^2\|w-w_0\|_2^2 .
\]
Since \(P(w)\asymp\rho\),
\[
  \Eb_{w_0}\{\ell_n(w_0)-\ell_n(w)\}
  \ge
  c\gamma n^2\rho\,\|w-w_0\|_2^2 .
\]

We next prove consistency.  With probability tending to one,
\[
  \sum_{e\in\mathcal D_n}A_e\le Cn^2\rho
\]
by Bernstein's inequality.  On this event, for all \(w,w'\in\calW\),
\[
  |\ell_n(w)-\ell_n(w')|
  \le Cn^2\rho\,\|w-w'\|_2 .
\]
Indeed, on an edge \(A_e=1\),
\(|\partial_w\log P_e(w)|=O(\rho)/O(\rho)=O(1)\), while on a non-edge
\(A_e=0\),
\(|\partial_w\log\{1-P_e(w)\}|=O(\rho)\).  Hence the sum of edge
contributions is \(O(n^2\rho)\) on the displayed event and the sum of
non-edge contributions is also \(O(n^2\rho)\).

Fix \(\varepsilon>0\).  Cover
\(\{w:\|w-w_0\|\ge\varepsilon\}\) by an \(\eta\)-net with
\[
  \eta=c\gamma\varepsilon^2
\]
and cardinality at most \((C/\eta)^K\).  The Lipschitz error is at most
half the expected separation.  For a fixed net point \(w\), the
log-likelihood ratio
\[
  \ell_n(w)-\ell_n(w_0)
\]
is a sum of independent bounded terms.  The probability band implies
\[
  \left|
  \log\frac{P_e(w)}{P_e(w_0)}
  \right|\le C,
  \qquad
  \left|
  \log\frac{1-P_e(w)}{1-P_e(w_0)}
  \right|
  \le C\rho ,
\]
and its variance is \(O(n^2\rho)\).  Since its mean is at most
\(-c\gamma n^2\rho\varepsilon^2\), Bernstein's inequality gives
\[
  \Pb_{w_0}
  \left\{
    \ell_n(w)-\ell_n(w_0)
    >
    -c\gamma n^2\rho\varepsilon^2/2
  \right\}
  \le
  \exp\{-c\gamma^2 n^2\rho\varepsilon^4\}.
\]
Because \(K\) is fixed and \(\gamma^3n^2\rho\to\infty\), the net union
probability tends to zero.  Thus
\[
  \sup_{\|w-w_0\|\ge\varepsilon}\ell_n(w)<\ell_n(w_0)
\]
with probability tending to one, proving \(\widehat w\to_p w_0\).

At \(w_0\),
\[
  \dot M(w_0)^\top\dot M(w_0)
  =
  M^\top M+O(n^2\rho^3),
\]
so \(\rho=o(\gamma)\) gives
\[
  \lambda_{\min}\{\dot M(w_0)^\top\dot M(w_0)\}
  \ge
  c\gamma n^2\rho^2 .
\]
Since \(P_{ij}(w_0)\asymp\rho\),
\[
  \lambda_{\min}\{\calI_n(w_0)\}
  \ge c\gamma n^2\rho .
\]
The score has mean zero and covariance \(\calI_n(w_0)\).  Moreover,
\[
  \max_{i<j}
  \frac{
    \dot P_{ij}(w_0)^\top\calI_n(w_0)^{-1}\dot P_{ij}(w_0)
  }{
    P_{ij}(w_0)\{1-P_{ij}(w_0)\}
  }
  \le
  \frac{C\rho^2}{(\gamma n^2\rho)\rho}
  =
  O((\gamma n^2)^{-1})
  \to0.
\]
Lindeberg's theorem gives
\[
  \calI_n(w_0)^{-1/2}S_n(w_0)\Rightarrow N(0,I_K).
\]

It remains to control the Hessian and Taylor remainder.  A direct
differentiation of the Bernoulli score gives
\[
  -\nabla^2\ell_n(w_0)-\calI_n(w_0)
  =
  \sum_{e\in\mathcal D_n}(A_e-P_e(w_0))\,B_e(w_0),
\]
where each \(B_e(w_0)\) is a \(K\times K\) matrix satisfying
\[
  \|B_e(w_0)\|_{\op}\le C.
\]
Indeed,
\[
  \frac{\ddot P_{e,kl}}{P_e(1-P_e)}=O(\rho),
  \qquad
  \frac{\dot P_{e,k}\dot P_{e,l}}{P_e^2(1-P_e)^2}=O(1).
\]
Hence for every entry,
\[
  \Var\!\left\{\sum_e(A_e-P_e)B_{e,kl}\right\}
  \le
  C\sum_eP_e(1-P_e)
  \le
  Cn^2\rho.
\]
Since \(K\) is fixed, entrywise and operator norms are equivalent up to
constants, and therefore
\[
  \|\nabla^2\ell_n(w_0)+\calI_n(w_0)\|_{\op}
  =
  O_p(\sqrt{n^2\rho})
  =
  o_p(\gamma n^2\rho),
\]
because \(\gamma^2n^2\rho\to\infty\), which follows from
\(\gamma^3n^2\rho\to\infty\) and \(\gamma\le1\).

The third derivative tensor has the analogous decomposition: each entry
is a deterministic sum of \(O(\rho)\) terms over \(\binom{n}{2}\) dyads plus a
centered sum \(\sum_e(A_e-P_e)C_e(w)\) with \(|C_e(w)|\le C\).  Thus the
deterministic part is \(O(n^2\rho)\), and the centered part is
\(O_p(\sqrt{n^2\rho})\).  Hence
\[
  \sup_{w\in\calW}\|\nabla^3\ell_n(w)\|_{\op}
  =
  O_p(n^2\rho).
\]
On Fisher balls of radius \(O(1)\), the Euclidean radius is
\(O\{(\gamma n^2\rho)^{-1/2}\}\).  The Fisher-normalized score remainder
is therefore
\[
  O_p(n^2\rho)\,(\gamma n^2\rho)^{-1}
  (\gamma n^2\rho)^{-1/2}
  =
  O_p\{(\gamma^3n^2\rho)^{-1/2}\}
  =
  o_p(1).
\]
Taylor expansion of the score gives
\[
  0=S_n(\widehat w)
  =
  S_n(w_0)-\calI_n(w_0)(\widehat w-w_0)+R_n,
  \qquad
  \|\calI_n(w_0)^{-1/2}R_n\|=o_p(1).
\]
Solving proves the linear expansion, rate, and CLT.
\end{proof}

\end{document}
